\documentclass[aspectratio=169, 14pt]{beamer}
\usepackage{beamer_preamble}

\begin{document}

% ==== Title ===============================================================
\titleslide{Lesson 5-4 \textperiodcentered\ Period 1}{Solving Radical Equations --- Foundation}

% ==== Do Now ==============================================================
\begin{frame}{Do Now --- Notice \& Wonder}{Algebra 2 \textperiodcentered\ Lesson 5-4 \textperiodcentered\ Period 1}
\phasebadges{2}{5}{notice \& wonder}
\vspace{0.2cm}
\begin{projcard}{Explore \& Reason \textperiodcentered\ Launch problem}{slidegold}
\textbf{A.} Solve \(3(a+1)^2 + 2 = 11\). Use at least two methods.

\vspace{0.15cm}
\textbf{B.} Try each method on \(3\sqrt{a+1} + 2 = 11\).

\vspace{0.15cm}
\textbf{C.} Which method works better with a radical? What problems arise with the other?
\end{projcard}
\end{frame}

% ==== Launch ==============================================================
\begin{frame}{Launch --- Isolate--Raise--Check Rules}{Algebra 2 \textperiodcentered\ Lesson 5-4 \textperiodcentered\ Period 1}
\phasebadges{1-2}{12}{teacher models}
\vspace{0.2cm}
\begin{projcard}{Isolate--Raise--Check Rules (keep on your page)}{slidegreen}
\begin{enumerate}
  \setlength{\itemsep}{0pt}
  \item \textbf{Isolate} the radical on one side of the equation.
  \item \textbf{Raise} both sides to the index power (\(\sqrt{\vphantom{x}}\) \(\rightarrow\) square; \(\sqrt[3]{\vphantom{x}}\) \(\rightarrow\) cube).
  \item Solve the resulting polynomial equation.
  \item \textbf{Check} each candidate in the \emph{original} equation --- reject any extraneous solutions.
\end{enumerate}

\vspace{0.1cm}
\textcolor{slideaccent}{\textbf{COMMON ERROR:} Squaring is NOT reversible --- always substitute into the ORIGINAL.}
\end{projcard}
\end{frame}

% ==== Explore =============================================================
\begin{frame}{Explore --- Think-Pair-Share (33 min)}{Algebra 2 \textperiodcentered\ Lesson 5-4 \textperiodcentered\ Period 1}
\phasebadges{1-2}{33}{student work}
\small\textcolor{slidegray}{6 items in order. Plan \(\sim\)5--6 min per item. Cut \#4 if pace slows.}

\vspace{0.1cm}
\begin{projcard}{Try It 1 \textperiodcentered\ Isolate and raise to solve}{slideblue}
\(\sqrt{x-2}+3=5\) \quad and \quad \(\sqrt[3]{x-1}=2\)
\end{projcard}
\vspace{-0.15cm}
\begin{projcard}{Try It 3 \textperiodcentered\ Solve + check for extraneous}{slideblue}
\(x = \sqrt{7x+8}\) \quad and \quad \(x+2 = \sqrt{x+2}\)
\end{projcard}
\vspace{-0.15cm}
\begin{projcard}{Practice \#22 + \#8 \textperiodcentered\ Solve radical equations}{slideblue}
\(\sqrt{4x}=11\) \quad then \quad \(\sqrt{8x+9}=x\) \quad --- substitute into ORIGINAL to check
\end{projcard}
\vspace{-0.15cm}
\begin{projcard}{Practice \#15 \textperiodcentered\ Identify extraneous (generalize)}{slideblue}
Explain how to identify an extraneous solution for a radical equation.
\end{projcard}
\vspace{-0.15cm}
\begin{projcard}{Practice \#4 \textperiodcentered\ Error Analysis (extension)}{slideblue}
Nazim says \(-3\) and \(6\) solve \(\sqrt{3x+18}=x\). What error did he make?
\end{projcard}
\end{frame}

% ==== Share / Summary =====================================================
\begin{frame}{Share / Summary}{Algebra 2 \textperiodcentered\ Lesson 5-4 \textperiodcentered\ Period 1}
\phasebadges{2}{5}{academic conversation}
\vspace{0.2cm}
\begin{projcard}{Sentence frame \textperiodcentered\ Pair share}{slideblue}
``I isolated the radical to get \underline{\hspace{1.5cm}}.
Squaring both sides gives \underline{\hspace{1.5cm}}.
My candidate solutions are \underline{\hspace{1.5cm}}.
When I check in the ORIGINAL equation, \underline{\hspace{1.5cm}} is valid
and \underline{\hspace{1.5cm}} is extraneous because \underline{\hspace{1.5cm}}.''
\end{projcard}

\vspace{0.15cm}
\begin{projcard}{Bridge to Period 2}{slideaccent}
Next class: equations with two radicals or a radical and a binomial.\\
Same Isolate--Raise--Check logic --- more complex setup.
\end{projcard}
\end{frame}

% ==== Exit Ticket =========================================================
\begin{frame}{Exit Ticket --- Summary Recap}{Algebra 2 \textperiodcentered\ Lesson 5-4 \textperiodcentered\ Period 1}
\phasebadges{1-2}{2}{not graded}
\vspace{0.2cm}
\begin{projcard}{Complete on your exit ticket}{slidegold}
\begin{enumerate}
  \setlength{\itemsep}{0.2cm}
  \item To solve a radical equation I first \underline{\hspace{5cm}}.
  \item I raise both sides to power \underline{\hspace{2cm}} because \underline{\hspace{3.5cm}}.
  \item An extraneous solution happens when \underline{\hspace{4cm}}.
\end{enumerate}
\end{projcard}
\end{frame}

\end{document}
